\section{Equações Lineares Elípticas de 2ª ordem}



\begin{frame}{Equações Elípticas de Segunda Ordem}
Desejamos estudar o problema de valor de fronteira
\[
\begin{cases}
{\color{red}L}{\color{blue}u}=f, & \text{ em } \Omega;\\
{\color{blue}u}=0, & \text{ sobre } \partial \Omega,
\end{cases}
\]
onde \(\Omega\in \rd\) é aberto e limitado e 
\begin{itemize}
\item {\color{blue}\(u:\overline{\Omega}\longrightarrow \R\)}
é uma função desconhecida;
\item \(f:\Omega \longrightarrow \R\) é uma função dada;
\item {\color{red}\(L\)} é um operador diferencial parcial de 2ª ordem que possui uma das seguintes formas:
\end{itemize}
%{\color{red}
%\[
%Lu=-\sum_{i,j=1}^{\dm}\left(a^{ij}(x)u_{x_i}\right)_{x_j}+\sum_{i=1}^{\dm}b^i(x)u_{x_i}+c(x)u
%\]
%}
%ou 
%{\color{red}
%\[
%Lu=-\sum_{i,j=1}^{\dm}a^{ij}(x)u_{x_ix_j}+\sum_{i=1}^{\dm}b^i(x)u_{x_i}+c(x)u
%\]
%}
\end{frame}
%%%%%%%%%%%%%%%%%%%%%%%%%%%%%%%%%%%%%%%%%%%%%%%%%%

\begin{frame}{Operador Elíptico}
\begin{itemize}
\item Forma Divergente
{\color{red}
\[
Lu=-\sum_{i,j=1}^{\dm}\left(a^{ij}(x)u_{x_i}\right)_{x_j}+\sum_{i=1}^{\dm}b^i(x)u_{x_i}+c(x)u
\]
}
\item Forma não divergente
{\color{red}
\[
Lu=-\sum_{i,j=1}^{\dm}a^{ij}(x)u_{x_ix_j}+\sum_{i=1}^{\dm}b^i(x)u_{x_i}+c(x)u
\]
}
\end{itemize}
onde \(a^{ij},b^i, c :\Omega \longrightarrow \R\) são funções dadas com \(i,j\in \ind\).
\end{frame}
%%%%%%%%%%%%%%%%%%%%%%%%%%%%%%%%%%%%%%%%%%%%%%%%%%
\subsection*{Definição de Soluções Fracas}
\begin{frame}{Soluções Fracas}
\begin{defin}
Dados \(a^{ij},b^i, c\in L^\infty(\Omega)\), para quaisquer \(i,j\in \ind\), definiremos a forma bilinear {\color{blue}\(B:\ho\times \ho\longrightarrow \R\) associada ao operador L na forma divergente}, por 
{\color{red}\[
B(u,v)=\int_\Omega \sij a^{ij}u_{x_i}v_{x_j}+\si b^iu_{x_i}v+cuv\,dx.
\]}
\end{defin}


\end{frame}
%%%%%%%%%%%%%%%%%%%%%%%%%%%%%%%%%%%%%%%%%%%%%%%%%%
\begin{frame}{Soluções Fracas}
\begin{defin}
Sejam{\color{red} \(f\in H^{-1}(\Omega)\)} e {\color{red}\(B\)} a forma bilinear associada ao operador {\color{red} \(L\)} na forma divergente. 
Dizemos que {\color{blue}\(u\in \ho\)} é uma {\color{blue}solução fraca} do problema 
{\color{red}\[
\begin{cases}
{L}{\color{blue}u}=f, & \text{ em } \Omega;\\
{\color{blue}u}=0, & \text{ sobre } \partial \Omega,
\end{cases}
\]}
se 
{\color{red}
\[B({\color{blue}u},{\color{black}v})=\innp{f,{\color{black}v}}_{H^{-1},H_0^1},\ \forall\, {\color{black} v}\in \ho.\]
}
\end{defin}
\end{frame}
%%%%%%%%%%%%%%%%%%%%%%%%%%%%%%%%%%%%%%%%%%%%%%%%%%
\subsection*{Teorema de Lax-Milgram}
\begin{frame}{Teorema de Lax-Milgram}
\begin{teo}[Lax-Milgram]
Seja \(H\) um {\color{blue}espaço de Hilbert real}, e suponhamos que \(B:H\times H\longrightarrow \R\) seja uma {\color{blue}forma bilinear contínua}. Suponhamos também que exista {\color{blue}\(\beta>0\)} tal que
\[
B(u,v)\geq {\color{blue}\beta} \n{u}^2 \text{ para todo } u\in H. \quad {\color{red}\text{(coercividade)}}
\]
 Então, para cada \(\varphi\in H'\), existe um único {\color{red}\(u_0\in H\)} tal que 
\[
B({\color{red}u_0},v)=\innp{\varphi,v}_{H',H},\ \text{ para todo } v\in H.
\]
\end{teo}
\end{frame}
%%%%%%%%%%%%%%%%%%%%%%%%%%%%%%%%%%%%%%%%%%%%%%%%%%
\subsection*{Elipticidade}
\begin{frame}{Elipticidade}
\begin{defin}
Dizemos que o operador diferencial {\color{red}\(L:\Omega\longrightarrow \R \)} é {\color{blue}uniformemente elíptico} se exister \(\theta>0\) tal que
{\color{red}\[
\sij a^{ij}(x)\xi_i\xi_j\geq \theta |\xi|^2,
\]}
para quase todo \(x\in \Omega\) e todo \(\xi=(\xi_1,\ldots,x_\dm)\in \rd\). 
\end{defin}

\end{frame}

%%%%%%%%%%%%%%%%%%%%%%%%%%%%%%%%%%%%%%%%%%%%%%%%%%

\begin{frame}
\begin{obs}
\begin{enumerate}[a]
\item Daqui em diante, assumiremos a {\color{red}condição de simetria}, isto é, 
\[{\color{red}a^{ij}=a^{ji},\ \forall\, i,j\in \ind.}\]
\item \(Lu=-\Delta u\) é um operador elíptico.
\item Se {\color{red}\(L\) é uniformemente elíptico}, então, para cada \(x\in \Omega\), a matriz
\[
A(x)=[a^{ij}(x)]_{i,j=1}^\dm
\]
é {\color{red}positiva e definida}. 
\end{enumerate}
\end{obs}
\end{frame}